\section{The ballistic element at terahertz frequencies}
\label{sec:ballistic}

One mechanism is outside the fluid model by construction. In the clean
limit the ballistic funnel of Sec.~\ref{sec:dc} has a curvature
$G_{2}=3.8$~mA/V$^{2}$; if it responded instantaneously, a cosine drive of
$0.1$~V across the device would rectify $18\ \mu$A, which is $0.87\AW$ of
the available power of a $75\ \Omega$ source, two to ten times the plasmonic
figures. Whether it does respond is a transit-time question. It has been
asked before of a diode of about $100\nm$, with a time-dependent particle
simulation that found rectification into the terahertz range and none
above about $5\THz$~\cite{stearns2020high}; we ask it of the device with
the documented dc asymmetry, then of its length, and for the shortest
device we compute the field as well.

The imposed potential is made time dependent, $U(x,t)=-eV_{a}\cos(\omega
t)\,g(x)$, and the flights integrate $\hbar\dot{\bm{k}}=-\nabla U$ without
the energy conservation that holds only in a static landscape. Two profiles
$g$ bracket the electrostatics, which the model does not compute: the
uniform ramp over the $6\um$ device used at dc, and a drop of length
$\ell_{F}$ centred on the neck, with $\ell_{F}=175$, $500$ and $1600\nm$
(one neck width, three, and the taper length). The bias-window estimator of
Eq.~\eqref{eq:landauer} relies on flux conservation in a fixed landscape and
fails in a moving one, so electrons are injected from both contacts with the
equilibrium distribution $N(E)f(E-\EF)$ at uniformly distributed phases, and
\begin{equation}
\Idc=\frac{g e}{h}\,\Phi\,\bigl[\langle T_{SD}\rangle-\langle T_{DS}\rangle\bigr],
\label{eq:twosided}
\end{equation}
with $\Phi=\int dE\,N(E)f(E-\EF)$ and the averages taken over injection
phase; $(ge/h)\Phi$ is $1.43$~mA
for the $2\um$ contacts. The first harmonic of the drain current, from the
arrival phases, gives the transport admittance and hence the mismatch to
the antenna. Six further tests validate the estimator: it reproduces the
bias-window current in a static field, it gives zero at zero bias, zero for
the straight control at every frequency, and at $300$~MHz it reproduces the
rectified current and the admittance computed from the static curve
(Supplemental Material). The results are reported as the transfer
$S_{B}(f)=\Idc(f)/\Idc(0)$, with $\Idc(0)$ the quasi-static value for the
same field profile, and as responsivity in two forms. $\RI$ is the current
of one device per available power of the $75\ \Omega$ source, with the
mismatch taken from the first-harmonic admittance. The ballistic devices
have resistances of $0.6$ to $4$~k$\Omega$, far from the antenna, so we
also give the current per watt absorbed at the static conductance $G$,
\begin{equation}
R_{\mathrm{abs}}=\frac{2\,\Idc}{V_{a}^{2}G},
\label{eq:rabs}
\end{equation}
which is what $N=1/(\Rant G)$ devices in parallel would give on a matched
antenna if their impedance did not change with frequency. The plasmonic
funnel of Sec.~\ref{sec:plasmonic} absorbs $87$ to $100\%$ of the available
power above $0.1\THz$, so its $\RI$ is to be compared with
$R_{\mathrm{abs}}$. We do not normalise by the power that the model itself
absorbs at the drive frequency: with an imposed field and no displacement
current its conduction admittance falls, vanishes and changes sign with
frequency, and a responsivity so defined diverges where it does. For
the one device whose field is computed (Sec.~\ref{sec:pic}) the absorbed
power is a result, and $R_{\mathrm{abs}}$ there is the rectified current
over it.

\subsection{The $6\um$ funnel in a uniform field: two roll-offs}

With the uniform field the response dies in two steps
[Fig.~\ref{fig:ballistic-thz}(a)]. The transfer is $0.98\pm0.04$ at
$0.3\GHz$ and already $0.82$ at $3\GHz$ and $0.67$ at $10\GHz$, although a
direct flight from the source contact to the neck takes $3.6$~ps. The
reason is in the distribution of flight times
[Fig.~\ref{fig:ballistic-thz}(c)]. The taper scatters part of the incoming
electrons into directions nearly perpendicular to the axis, and these drift
back to the source contact, bouncing between the walls of the lead, over
tens of picoseconds: $5\%$ of the source-injected electrons fly longer than
$20$~ps and $0.9\%$ longer than $50$~ps. In a field that fills the leads,
this slow population is displaced in energy between wall collisions once
the period approaches its lifetime, and the transmission in both directions
changes. The response time of a ballistic device is set by its slowest
population, not by its mean transit, and a converging taper manufactures
slow electrons. The second step is the transit roll-off of the direct
electrons: the transfer is $0.41$ at $100\GHz$, crosses zero at $300\GHz$,
where $\sin(\pi f\tau_{\mathrm{tr}})$ vanishes for
$\tau_{\mathrm{tr}}=3.6$~ps, and is $0.05\pm0.04$ at $1\THz$, i.e.
$0.05\AW$. At a mean free path of $700\nm$ the quasi-static responsivity is
only $0.15\AW$ to begin with; it is $0.04\AW$ at $30\GHz$ and zero at
$1\THz$. With the field
model in which the dc asymmetry was established, the $6\um$ funnel is not
a terahertz rectifier.

\subsection{Field confined to the neck}

Self-consistent electrostatics concentrates the field of such a diode at
the neck rather than spreading it along the
leads~\cite{truccolo2024comprehensive}. Confining the drop to the
neck lowers the static asymmetry, from $1.61$ for the uniform field to
$1.52$, $1.50$ and $1.40$ for $\ell_{F}=1600$, $500$ and $175\nm$, and it
changes the dynamics [Fig.~\ref{fig:ballistic-thz}(a) and
\ref{fig:ballistic-thz}(b)]. For $\ell_{F}=500\nm$ the transfer is
$1.02\pm0.04$ up to $3\GHz$, since the slow electrons in the lead no
longer see a field, then falls to $0.54$ at $100\GHz$, and is $0.64$ at
$300\GHz$ and $0.38\pm0.04$ at $1\THz$. The decline between $10$ and
$100\GHz$ is not what the $0.5$~ps crossing of the field region would
suggest; the transfer is set by the whole flight between the contacts. At
$1\THz$ the transfer is $0.51$ for a drop of one neck width, $0.38$ for
$500\nm$ and zero for $1600\nm$. The corresponding responsivities of one device are
$0.45\pm0.04$ and $0.37\pm0.04\AW$, as much as the Dyakonov-Shur response
of the plasmonic funnel at $100$~K; per absorbed watt the $500\nm$ drop
gives $0.72\pm0.07\AW$, about twice that. The
amplitude dependence is that of the dc curve: at $300\GHz$ the rectified
current follows $V_{a}^{2}$ to $0.1$~V and is $0.40$ and $0.24$ of the
square law at $0.2$ and $0.3$~V.

Two further cases show that the terahertz response of the confined field is
not simply the static asymmetry carried up in frequency. At $\lambda=700\nm$
the static asymmetry of the confined-field models is reversed ($A=0.94$ to
$0.97$), the rectified current passes through zero between $30$ and
$100\GHz$, and it returns with the sign of the clean funnel, $0.20\AW$ at
$300\GHz$ and $0.24\AW$ at $1\THz$ ($0.66$ and $0.75\AW$ per absorbed
watt). With half-diffuse walls ($p=0.5$) in the
clean limit the static asymmetry of the $500\nm$ drop is $1.01$, and the
responsivity at $1\THz$ is nevertheless $0.45\pm0.04\AW$ ($0.31\AW$ for the
neck-width drop). There is, then, a transit-time rectification of the
funnel with a neck-confined field that is robust against edge roughness
and bulk scattering, of $0.2$ to $0.45\AW$ for one device. It does not
depend on the treatment of the few electrons that reach the Dirac point:
replacing their removal by a classical turning point changes no result by
more than one standard error. It does depend on the confinement of the
field, which this model assumes and cannot test: at $1\THz$ and
$\lambda=700\nm$ the $6\um$ device gives $0.00\pm0.04\AW$ per absorbed
watt with the uniform field and $0.75\pm0.05\AW$ with the $500\nm$ drop.
For this device the terahertz response is decided by electrostatics that
the model does not contain.

\begin{figure*}
\includegraphics[width=\textwidth]{\figpath/fig7_ballistic_thz.pdf}
\caption{The $6\um$ ballistic funnel under an ac drive of amplitude $0.1$~V. (a)
Rectified current relative to its quasi-static value for the uniform field
(clean limit and $\lambda=700\nm$) and for a potential drop of $500\nm$ at
the neck; the dotted line is the transit factor for a flight of $3.6$~ps
from the source contact to the neck. (b) The same ratio at $0.3$ and
$1\THz$ against the length of the drop; open symbols: uniform field. (c)
Fraction of electrons with a flight time longer than $t$ in the static
device at $0.1$~V. Error bars are one standard error.}
\label{fig:ballistic-thz}
\end{figure*}

\subsection{Length scaling}
\label{sec:length}

The two roll-offs of the $6\um$ device are properties of its size. The
model has no length of its own (Sec.~\ref{sec:model-ballistic}): a copy of
the device at half scale has half the current at every bias, and under a
drive at twice the frequency it has half the rectified current. Both
statements were tested, as was the equivalence with scattering at fixed
$\lambda/L$: the static current of device D at $0.1$~V and twice that of its
half-scale copy are $188.0$ and $188.6\ \mu$A, and the rectified current
at $10\GHz$ and twice that of the copy at $20\GHz$ are $11.9\pm0.4$ and
$12.1\pm0.5\ \mu$A. The
frequency response of a given shape is therefore a function of the reduced
frequency $fL/\vF$, and the question whether a ballistic funnel rectifies
at $1\THz$ becomes the question how long it is.

Table~\ref{tab:short} and Fig.~\ref{fig:length}(a) define five devices: the
$6\um$ funnel (G0), its half-scale copy (G1), the same taper and neck with
the leads cut to $250\nm$ (G2), a $0.55\um$ funnel with a $100\nm$ neck
(G3), and a $0.50\um$ funnel with the shoulder and neck of the diode
measured at $28\THz$ in Ref.~\cite{zhu2013graphene}, for which we assume a
45 degree taper (G4). Mirror-symmetric versions of G2 to G4 serve as
controls; at $0.1$, $1$ and $10\THz$ and for both field profiles all 22
control points are zero within $2.1$ standard errors. The confined drop is
one neck width long in every device.

With a uniform field and no scattering the five transfer functions fall on
one curve [Fig.~\ref{fig:length}(b)]: above $fL/\vF=0.2$ they agree within
about $0.2$, the response is half its quasi-static value near
$fL/\vF=0.4$, and it is zero from about $1.7$ on; for device D it remains
zero within $0.04$ out to $fL/\vF=30$. The frequency of half response is
thus $0.4\,\vF/L$: $60\GHz$ for the $6\um$ device and $0.7\THz$ for the
$0.55\um$ one, and the response is gone at $280\GHz$ and at $3\THz$. The
same scaling with the size of the scatterer over the Fermi velocity was
found for the ratchet current of Dirac electrons on asymmetric
antidots~\cite{ermann2011relativistic}. Below $fL/\vF=0.2$ the curves
differ, and the difference identifies the gigahertz step of the $6\um$
device. Devices G0 and G1, with leads as long as the taper, share a
plateau at $0.67$ between $fL/\vF=0.03$ and $0.3$; G2 and G3, with short
leads, are still at $0.94$ and $0.90$ at $0.06$, and G4 lies between. The slow population
belongs to the leads, as inferred above from the flight times. Scattering
lowers the cut-off when the device is longer than the mean free path and
not otherwise: for G3 at $\lambda=700\nm$ ($\lambda/L=1.3$) the transfer
is that of the clean device, $0.76$ and $0.32$ at $0.3$ and $1\THz$,
while for G2 ($\lambda/L=0.34$) it is $0.41$ against $0.75$ at
$fL/\vF=0.2$. A cut-off near the Drude damping rate, as reported for the
simulated diode of Ref.~\cite{stearns2020high}, and a cut-off at the
flight time are the two limits of the same behaviour.

\begin{table*}
\caption{Funnels used for the length scaling. Shoulder width, neck width
and lead lengths (source, drain) in nm; $L$ is the length between the
contacts. $A$ is the static asymmetry at $0.1$~V with a uniform field for
no scattering and for mean free paths of $700$ and $100\nm$; $R$ is the
resistance at $\lambda=700\nm$ and $N=1/(\Rant G)$ the number of devices in
parallel that matches $75\ \Omega$. G4 has the shoulder and neck of the
diode of Ref.~\cite{zhu2013graphene} with an assumed 45 degree taper; the
others have the 30 degree taper of device D.}
\label{tab:short}
\begin{ruledtabular}
\begin{tabular}{llcccccccc}
 & Description & Shoulder/neck & Leads & $L$ ($\mu$m) & $A$, clean & $A$, 700 nm & $A$, 100 nm & $R$ (k$\Omega$) & $N$ \\
\hline
G0 & device D of Ref.~\cite{nguyen2021geometric} & 2000/175 & 2000, 2420 & 6.00 & 1.61 & 1.10 & & 0.95 & 13 \\
G1 & G0 at half scale & 1000/88 & 1000, 1210 & 3.00 & 1.61 & 1.13 & 1.06 & 1.50 & 20 \\
G2 & G0 with short leads & 2000/175 & 250, 250 & 2.08 & 1.91 & 1.55 & 1.33 & 0.77 & 10 \\
G3 & short funnel & 500/100 & 100, 100 & 0.55 & 1.83 & 1.56 & 1.28 & 1.16 & 15 \\
G4 & proportions of Ref.~\cite{zhu2013graphene} & 400/75 & 169, 169 & 0.50 & 1.60 & 1.30 & 1.15 & 1.54 & 21 \\
\end{tabular}
\end{ruledtabular}
\end{table*}

\subsection{A short funnel at $1\THz$ with an imposed field}
\label{sec:short}

For the $0.55\um$ funnel the reduced frequency at $1\THz$ is $0.55$, and
the device is shorter than the room-temperature mean free path. At
$\lambda=700\nm$ and $1\THz$ the imposed-field model gives
\begin{equation}
R_{\mathrm{abs}}=
\begin{cases}
0.70\pm0.02\AW & \text{uniform field},\\
0.71\pm0.03\AW & \text{drop at the neck},
\end{cases}
\label{eq:imposed}
\end{equation}
and $\RI=0.17$ and $0.14\AW$ for one device of $1.2$~k$\Omega$ on the
$75\ \Omega$ antenna [Fig.~\ref{fig:length}(c)]; G4 gives $0.57\pm0.03$ and
$0.70\pm0.02\AW$ per absorbed watt. Two things have changed with respect to
the $6\um$ device. The two profiles, which differed by the whole effect at
$6\um$, agree within the statistical error [Fig.~\ref{fig:length}(d)]. And
the result holds at the mean free path of encapsulated graphene at room
temperature, not only in the clean limit, in which the same device gives
$0.95\pm0.05$ and $0.83\pm0.06\AW$. Taken at face value this is two to
three times the response of the plasmonic funnel at $100$~K and eight
times its $0.09\AW$ at $300$~K.

In the same model the edges and scattering cost a factor of up to five.
With half-diffuse walls ($p=0.5$) the figures at $1\THz$ and
$\lambda=700\nm$ are $0.42\AW$ (uniform) and $0.14\AW$ (confined); with a
mean free path of $100\nm$, which is also a rough bound for the effect of
electron-electron collisions at $300$~K, they are $0.22\pm0.04$ and
$0.42\pm0.04\AW$. With the confined drop the static asymmetry of G3 falls
below one when scattering is strong ($A=0.81$ at $100\nm$, $0.90$ with
half-diffuse walls at $700\nm$), while at $1\THz$ the current has the sign
of the clean funnel in every case, as for the $6\um$ device. The measured
devices bear the dependence on the mean free path out in order of
magnitude: the diode of Ref.~\cite{zhu2013graphene}, with a measured mean
free path of about $45\nm$, had a zero-bias dc responsivity of $0.12\AW$,
against $0.71\AW$ for G4 at $100\nm$ and $1.33\AW$ at $700\nm$ in the
quasi-static limit of the uniform-field model, and the rectifiers of
chemical vapor deposited graphene of Ref.~\cite{hemmetter2026enhanced},
with mean free paths of at most $65\nm$, have a responsivity of $0.8$~V/W
referred to the beam power, against hundreds of V/W for encapsulated
ballistic rectifiers~\cite{auton2017terahertz}. The amplitude dependence
is a square law with a low ceiling: with the uniform field the rectified
current of G3 follows $V_{a}^{2}$ up to $0.1$~V and is $0.80$ and $0.57$ of
the law at $0.2$ and $0.3$~V, and with the confined drop it is already
$19\%$ below the law at $0.1$~V. The scale is the Fermi energy,
$\EF/e=0.14$~V, a tenth of the gate overdrive that limits the plasmonic
element.

Above $1\THz$ the two field profiles separate again. With the uniform field
the response of G3 and G4 is zero at $3\THz$, as the scaling requires. With
the drop confined to the neck it is $0.85\pm0.03$ and $1.03\pm0.03\AW$ at
$3\THz$ and has fallen by an order of magnitude at $10\THz$, to
$0.05\pm0.03$ and $0.09\pm0.03\AW$ at $\lambda=700\nm$ ($0.23\pm0.04\AW$
for G4 at $100\nm$); at $30\THz$ no device, profile or mean free path gives
a current distinguishable from zero. The response at $28\THz$ measured
for a diode of the size of G4~\cite{zhu2013graphene} is not reproduced. At
that frequency the photon energy, $116$~meV, is close to the Fermi energy
and a classical trajectory in a classical field is not an adequate
description; interband absorption, which the model lacks, is the natural
candidate for what was measured.

The agreement of the two profiles in Eq.~\eqref{eq:imposed} suggests that
for a device this short the field does not matter. The next subsection
tests that by computing it, and finds that it does.

\subsection{The short funnel with the field computed}
\label{sec:pic}

We take the $0.55\um$ funnel G3 into the model of
Sec.~\ref{sec:model-pic}, in which the potential follows from the charge
of the electrons, a gate and metal contacts. The operating point is that
of Eq.~\eqref{eq:imposed}: $\lambda=700\nm$, $1\THz$, an amplitude of
$0.1$~V between the contacts, the drive that leaves only the shape, and
the gate at $300\nm$. The funnel rectifies $0.429\pm0.012\ \mu$A of an
absorbed $3.44\ \mu$W,
\begin{equation}
R_{\mathrm{abs}}=\frac{\Idc}{P_{\mathrm{abs}}}=0.124\pm0.003\AW ,
\label{eq:computed}
\end{equation}
and $\RI=0.023\AW$ for one device on the $75\ \Omega$ antenna, with an
admittance of $0.69-0.17i$~mS. Under the $30\nm$ gate the figures are
$0.132\pm0.003$ and $0.026\AW$, with $0.73+0.34i$~mS. The absorbed power
is now a computed quantity; in the definition of Eq.~\eqref{eq:rabs}, for
comparison, the values are $0.10$ and $0.11\AW$ against $0.70\AW$. The
errors are statistical, from seven and six runs. Changes of the grid, the
time step, the particle number and the length of the contacts move the
current between $0.38$ and $0.49\ \mu$A, and the mirror-symmetric version
of the device, which must give zero under this drive, gives currents of up
to $0.06\ \mu$A in single runs at $0.3$ and $1\THz$ (the mean of six runs
from $0.3$ to $3\THz$ is $0.01\pm0.02\ \mu$A); we take $0.05\ \mu$A as the
floor of the method and quote $0.12\pm0.015\AW$. The computed field leaves a sixth of
what either imposed field gave.

The loss is already present at dc. The static asymmetry at $0.1$~V is
$1.053$ under both gates, against $1.56$ and $1.13$ for the two imposed
profiles, while the conductance, $0.86$~mS, is that of the uniform-field
model: only the even part of the current-voltage curve has changed.
Figure~\ref{fig:pic}(a) and \ref{fig:pic}(b) show why. About $61\%$ of the
bias drops within $90\nm$ of the neck, centred on it; $14\%$ along the
lead and the taper; and $25\%$ at the two contact planes, $6$ to $9\%$ at
the source and $17$ to $19\%$ at the drain, where it contributes nothing
to the asymmetry. The drop at the source is of the size expected for a
ballistic contact, half of the current over the Sharvin conductance of
the $500\nm$ lead, or $8\%$ of the bias. Neither imposed profile
resembles this. The uniform ramp puts most of the bias along the taper,
which is where its large asymmetry comes from, and the confined drop puts
all of it at the neck: with $61\%$ there and a rectified current that goes
as the square of the bias, the confined-drop model would give
$0.9\ \mu$A in the quasi-static limit, and the computed curve gives
$1.07\pm0.02\ \mu$A.

The frequency dependence is that of the uniform-field model at a seventh
to a ninth of its height [Fig.~\ref{fig:pic}(c)]. At $0.1\THz$ the
rectified current is $1.04\pm0.01\ \mu$A, the quasi-static value of the
computed static curve; it is $0.94$, $0.71$, $0.43$ and $0.15\ \mu$A at
$0.3$, $0.6$, $1$ and $1.8\THz$, half at $0.9\THz$ or $fL/\vF=0.5$. The
transit roll-off of Sec.~\ref{sec:length} thus survives the computation
of the field. The agreement of the two imposed profiles at $1\THz$ does
not: it was the crossing of a curve that falls with frequency and one
that rises, and no sign of insensitivity to the field. Above $1\THz$ the
gate matters. Under the $300\nm$ gate the absorbed power falls steadily,
to $0.8\ \mu$W at $3\THz$, with $0.19\pm0.02\ \mu$A rectified there and
nothing at $6\THz$. Under the $30\nm$ gate the device is a plasmonic
cavity between two contacts that hold the potential, with a half-wave
resonance near $3.6\THz$ and a full-wave one near $7\THz$ for the plasma
speed of $3.9\,\vF$; the antisymmetric shape-only drive couples to the
second, the absorbed power rises to $10.7\ \mu$W at $6\THz$, and the
rectified current stays at $0.35$ to $0.50\ \mu$A up to $3\THz$.

With a gate in the model, shape and circuit can be compared in the
ballistic gas as they were in the fluid [Figs.~\ref{fig:pic}(d) and
\ref{fig:pic}(e)]. The circuit term is the current of a straight strip
of the same length and lead width with its drain tied to the gate, the
Dyakonov-Shur termination, here without collisions. Under the $300\nm$
gate it is $0.13$, $0.34$ and $0.01\ \mu$A at $0.3$, $1$ and $3\THz$,
against $0.94$, $0.43$ and $0.19\ \mu$A for the shape term: the shape is
the larger at each frequency. Under the $30\nm$ gate the strip gives
$0.67$, $1.15$ and $-1.77\ \mu$A against $0.89$, $0.48$ and $0.50\ \mu$A.
At $1\THz$ the circuit term is then more than twice the shape term,
$0.20\AW$ per absorbed watt against $0.13\AW$, and the funnel with its
drain tied to the gate, which has both, gives $1.01\ \mu$A, no more than
the strip. At $3\THz$, near the half-wave resonance that this
termination does excite, the funnel absorbs four times its
low-frequency power and the strip's current is reversed; the run was
repeated with a longer settling time and did not change, and we have not
looked for the cause. With the gate closer than the device is long,
then, the gate connection rectifies more than the shape at $1\THz$ in
the ballistic gas as well.

How much of the shape term survives in a real device is again a matter of
scattering and of the edges, and here the computed field changes the
answer of the imposed one in both directions [Fig.~\ref{fig:pic}(f)].
The mean free path matters little. Per absorbed watt at $1\THz$ the
funnel gives $0.130\pm0.005\AW$ with no scattering, $0.124$ at $700\nm$
and $0.107\pm0.007\AW$ at $100\nm$ ($0.145$, $0.132$ and $0.099\AW$ under
the $30\nm$ gate), where the uniform-field model lost a factor of three.
Scattering does lower the static asymmetry, $1.09$, $1.05$ and $1.02$,
and the quasi-static current, $2.1$, $1.1$ and $0.23\ \mu$A, but the
current at $1\THz$ falls much less, $0.50$, $0.43$ and $0.22\ \mu$A: what
a long mean free path adds to the static asymmetry is not available at
$1\THz$, and at $100\nm$ the roll-off has moved above $1\THz$. The
conductance halves at $100\nm$, so one device on the antenna does lose a
factor of two, $\RI=0.012\AW$. The edges decide. With half-diffuse walls
($p=0.5$, $\lambda=700\nm$) the static asymmetry is $1.005$ and the
current at $1\THz$ is $-0.07\pm0.01\ \mu$A under the $300\nm$ gate and
$-0.13\pm0.01\ \mu$A under the $30\nm$ gate, a sixth to a quarter of the
specular current with the opposite sign; two independent runs agree, and
the mirror-symmetric control with the same walls gives zero within the
floor. The imposed fields had given $0.42$ and $0.14\AW$ for this case;
the computed field gives at most $0.04\AW$ in magnitude. The shape
rectification of the short funnel requires edges that reflect specularly.
Only $p=1$ and $p=0.5$ were computed, so how specular is not known.

The amplitude dependence differs from that of the imposed fields. At
$0.05$~V the current is $0.10\pm0.02\ \mu$A under both gates, consistent
with the square law through the $0.1$~V point. At $0.2$~V, which is above
$\EF/e$, it is $2.56$ and $2.66\ \mu$A, $1.5$ and $1.4$ times the square
law, where both imposed fields fell below it; the absorbed power follows
$V_{a}^{2}$ within $7\%$, so $R_{\mathrm{abs}}$ rises to $0.20\AW$. Into a
matched dc load the conversion efficiency estimated from the
short-circuit current and the static conductance,
$\eta=\Idc^{2}/(4GP_{\mathrm{abs}})$, is $2\times10^{-5}$ at $0.1$~V,
where the device absorbs $3.4\ \mu$W, and $1.5\times10^{-4}$ at $0.2$~V
and $13\ \mu$W. As a detector, $N=1/(\Rant\,\mathrm{Re}\,Y_{1})=19$
devices in parallel with their reactance tuned out absorb the available
power, and the Johnson noise of their static conductance gives a
noise-equivalent power of $130\pWrtHz$ at $1\THz$ and $300$~K, twice that
of the plasmonic funnel at the same temperature.

The outcome for the short funnel is therefore narrower than the imposed
fields suggested. A $0.55\um$ funnel with specular edges rectifies at
$1\THz$ by its shape, at $0.10$ to $0.13\AW$ per absorbed watt for any
mean free path from $100\nm$ up. That is the level of the plasmonic
funnel at $300$~K, $0.09\AW$, and below its $0.21$ to $0.34\AW$ at
$100$~K; it is not the factor of eight of Eq.~\eqref{eq:imposed}. With
rough edges it is lost, and under a close gate the circuit term exceeds
it.

\begin{figure*}
\includegraphics[width=\textwidth]{\figpath/fig10_ballistic_length.pdf}
\caption{Length scaling of the ballistic response at an ac amplitude of
$0.1$~V. (a) The five funnels of Table~\ref{tab:short} to scale, and the
two short ones enlarged. (b) Rectified current relative to its
quasi-static value against the reduced frequency, for a uniform field and
no scattering; the band marks the fall from half response to zero. (c)
Responsivity of the $0.55\um$ funnel G3 at $\lambda=700\nm$: per absorbed
watt, Eq.~\eqref{eq:rabs} (thick), and for one device on a $75\ \Omega$
antenna (thin, open symbols). (d) Responsivity per absorbed watt at
$1\THz$ against device length for the two field profiles. The blue band in
(c) and (d) is the range of the plasmonic funnel of
Sec.~\ref{sec:plasmonic}, $0.07$ to $0.34\AW$. Error bars are one standard
error.}
\label{fig:length}
\end{figure*}

\begin{figure*}
\includegraphics[width=\textwidth]{\figpath/fig11_selfconsistent.pdf}
\caption{The $0.55\um$ funnel G3 with the field computed; $\lambda=700\nm$
and an amplitude of $0.1$~V except in (f). (a) Static potential energy at
a bias of $0.1$~V with the gate at $300\nm$, in units of the bias;
contours every $0.1$; the dotted rectangles are the contact strips. (b)
The same along the axis for the two gate distances, against the two
imposed profiles of Sec.~\ref{sec:short}; colours and line styles as in
(c). (c) Rectified current against
frequency under the drive that leaves only the shape, for the imposed
fields and for the computed field with the gate at $300\nm$ (full
symbols) and $30\nm$ (open symbols); points within two standard errors of
zero are not drawn. (d), (e) Shape against circuit for the two gate
distances: the funnel under the shape-only drive, a straight strip of the
same length with its drain tied to the gate, the funnel with its drain
tied to the gate, and the mirror-symmetric funnel under the shape-only
drive, which must give zero. (f) Responsivity at $1\THz$ for no
scattering, mean free paths of $700$ and $100\nm$, and half-diffuse
walls: per absorbed watt for the computed field (diamonds, full and open
as in (c)), and by Eq.~\eqref{eq:rabs} for the imposed fields (circles:
uniform; squares: drop at the neck). The band is the range of the
plasmonic funnel. Error bars are one standard error.}
\label{fig:pic}
\end{figure*}
